\documentclass{article}
\usepackage[utf8]{inputenc}
\title{Pedestrian Crossing Behavior Review}
\begin{document}
\maketitle
\section{Introduction and Background}
This review synthesizes evidence from 12 papers examining Safety assessment, risk analysis, or safety performance evaluation, pedestrian crossing compliance analysis, pedestrian perception analysis, with findings organized around 55 discrete claims drawn from a structured comparison matrix. The reviewed evidence employs Empirical/Quantitative research, mixed methods, Microscopic simulation model / Behavioral model with gap acceptance as primary analytical frameworks. A persistent challenge in this literature is the heterogeneity in how studies operationalize crossing behavior constructs, report metrics, and handle context specificity — factors that collectively limit direct cross-study comparison. This review structures its analysis around methodological families, task-level performance patterns, metric reporting practices, and verified gaps (2 identified), moving beyond siloed paper summaries toward a structured evidence synthesis. The review is organized as follows: Section 2 examines methodological approaches and their representational tradeoffs; Section 3 addresses task-level findings and performance patterns; Section 4 evaluates metric reporting and comparability; Section 5 identifies verified research gaps and implications; Section 6 provides a field-level synthesis and research priorities.

\section{Research Methodologies in Pedestrian Crossing Studies}
The Empirical/Quantitative research family constitutes a primary methodological strand in the reviewed literature. Represents a distinct representational approach with characteristic strengths and limitations. Evidence across reviewed studies suggests that this approach is particularly well-suited to analyzing crossing decision determinants under specific traffic conditions. Limitations of this approach include reduced generalizability to contexts with substantially different traffic signal configurations. The mixed methods family represents a primary methodological strand in the reviewed literature. Represents a distinct representational approach with characteristic strengths and limitations. Evidence across reviewed studies suggests that this approach is particularly well-suited to analyzing crossing decision determinants under specific traffic conditions. Limitations of this approach include reduced generalizability to contexts with substantially different traffic signal configurations. The Microscopic simulation model / Behavioral model with gap acceptance family represents a secondary methodological strand in the reviewed literature. Represents a distinct representational approach with characteristic strengths and limitations. Evidence across reviewed studies suggests that this approach is particularly well-suited to analyzing crossing decision determinants under specific traffic conditions. Limitations of this approach include reduced generalizability to contexts with substantially different traffic signal configurations. The microscopic simulation social force model family represents a secondary methodological strand in the reviewed literature. Captures emergent group behaviors including collective waiting dynamics but requires careful parameter calibration. Evidence across reviewed studies suggests that this approach is particularly well-suited to analyzing pedestrian-vehicle interaction patterns. Limitations of this approach include reduced generalizability to contexts with substantially different group compositions. The microscopic traffic simulation agent based modeling family represents a secondary methodological strand in the reviewed literature. Represents a distinct representational approach with characteristic strengths and limitations. Evidence across reviewed studies suggests that this approach is particularly well-suited to analyzing crossing decision determinants under specific traffic conditions. Limitations of this approach include reduced generalizability to contexts with substantially different traffic signal configurations. Across these 12 methodological families, a clear tension emerges between interpretive richness and predictive performance. Structural models — including discrete choice and social force approaches — provide parameters directly interpretable for traffic engineering and policy purposes, whereas data-driven approaches achieve higher predictive accuracy on specific benchmark datasets but lack comparable policy utility. No single paradigm has been validated across the full range of pedestrian populations, traffic environments, and signal configurations represented in this corpus.

\section{Microscopic Simulation and Behavioral Modeling Approaches}
The Empirical/Quantitative research family constitutes a primary methodological strand in the reviewed literature. Represents a distinct representational approach with characteristic strengths and limitations. Evidence across reviewed studies suggests that this approach is particularly well-suited to analyzing crossing decision determinants under specific traffic conditions. Limitations of this approach include reduced generalizability to contexts with substantially different traffic signal configurations. The mixed methods family represents a primary methodological strand in the reviewed literature. Represents a distinct representational approach with characteristic strengths and limitations. Evidence across reviewed studies suggests that this approach is particularly well-suited to analyzing crossing decision determinants under specific traffic conditions. Limitations of this approach include reduced generalizability to contexts with substantially different traffic signal configurations. The Microscopic simulation model / Behavioral model with gap acceptance family represents a secondary methodological strand in the reviewed literature. Represents a distinct representational approach with characteristic strengths and limitations. Evidence across reviewed studies suggests that this approach is particularly well-suited to analyzing crossing decision determinants under specific traffic conditions. Limitations of this approach include reduced generalizability to contexts with substantially different traffic signal configurations. The microscopic simulation social force model family represents a secondary methodological strand in the reviewed literature. Captures emergent group behaviors including collective waiting dynamics but requires careful parameter calibration. Evidence across reviewed studies suggests that this approach is particularly well-suited to analyzing pedestrian-vehicle interaction patterns. Limitations of this approach include reduced generalizability to contexts with substantially different group compositions. The microscopic traffic simulation agent based modeling family represents a secondary methodological strand in the reviewed literature. Represents a distinct representational approach with characteristic strengths and limitations. Evidence across reviewed studies suggests that this approach is particularly well-suited to analyzing crossing decision determinants under specific traffic conditions. Limitations of this approach include reduced generalizability to contexts with substantially different traffic signal configurations. Across these 12 methodological families, a clear tension emerges between interpretive richness and predictive performance. Structural models — including discrete choice and social force approaches — provide parameters directly interpretable for traffic engineering and policy purposes, whereas data-driven approaches achieve higher predictive accuracy on specific benchmark datasets but lack comparable policy utility. No single paradigm has been validated across the full range of pedestrian populations, traffic environments, and signal configurations represented in this corpus.

\section{Psychological and Attitudinal Models of Signal Non-Compliance}
This section synthesizes findings from 12 papers examining Safety assessment, risk analysis, or safety performance evaluation, pedestrian crossing compliance analysis, pedestrian perception analysis using Empirical/Quantitative research, mixed methods, Microscopic simulation model / Behavioral model with gap acceptance. The comparative evidence matrix structures claims, metrics, and limitations to support structured analysis. A key observation is that Empirical/Quantitative research and mixed methods approaches are the dominant methodological families in this corpus, though significant heterogeneity in reporting practices limits direct cross-study comparison.

\section{Mobile Phone Distraction and Visual Attention}
This section synthesizes findings from 12 papers examining Safety assessment, risk analysis, or safety performance evaluation, pedestrian crossing compliance analysis, pedestrian perception analysis using Empirical/Quantitative research, mixed methods, Microscopic simulation model / Behavioral model with gap acceptance. The comparative evidence matrix structures claims, metrics, and limitations to support structured analysis. A key observation is that Empirical/Quantitative research and mixed methods approaches are the dominant methodological families in this corpus, though significant heterogeneity in reporting practices limits direct cross-study comparison.

\section{Geographic and Contextual Variation in Crossing Behavior}
This section synthesizes findings from 12 papers examining Safety assessment, risk analysis, or safety performance evaluation, pedestrian crossing compliance analysis, pedestrian perception analysis using Empirical/Quantitative research, mixed methods, Microscopic simulation model / Behavioral model with gap acceptance. The comparative evidence matrix structures claims, metrics, and limitations to support structured analysis. A key observation is that Empirical/Quantitative research and mixed methods approaches are the dominant methodological families in this corpus, though significant heterogeneity in reporting practices limits direct cross-study comparison.

\section{Gap Analysis and Synthesis of Findings}
The comparative analysis of 12 papers reveals 2 persistent gaps that current approaches have not adequately resolved. These gaps are documented through verified evidence from the review matrix and are categorized by severity below. Gap [4c163add] (medium severity): Current coverage appears to underrepresent Chinese-language or Chinese-context studies.. This gap limits the ability to draw robust conclusions from the existing literature. Gap [67b03b55] (medium severity): There are unresolved claim differences across papers under the same task, suggesting a need for clearer comparison criteria.. This gap limits the ability to draw robust conclusions from the existing literature. These gaps collectively indicate that the field lacks the empirical foundation needed to support confident generalization of crossing behavior models across contexts, populations, and technologies. Prioritizing the most severe gaps in future empirical work will be essential for advancing both scientific understanding and the practical applicability of crossing behavior models.

\section{Conclusions and Future Research Directions}
The evidence across 12 reviewed papers converges on a field-level picture in which pedestrian crossing behavior at signalized intersections is influenced by a complex interaction of individual characteristics, traffic flow properties, signal timing parameters, and environmental context. Methodological approaches to modeling this behavior range from structural models grounded in behavioral theory to data-driven approaches optimized for predictive accuracy. Three stable conclusions emerge from this review. First, waiting time is among the most consistent predictors of crossing initiation and violation behavior across reviewed studies. Second, methodological approaches exhibit a persistent tradeoff between interpretability and predictive performance. Third, metric reporting heterogeneity across studies substantially limits the precision of cross-study synthesis. Two persistent tensions run through this literature. The first concerns the cross-context validity of gap acceptance thresholds: values derived from studies in North American and European contexts may not transfer to other traffic environments with different vehicle-pedestrian interaction norms. The second tension concerns the behavioral effects of automated vehicles: whether pedestrians recalibrate their crossing heuristics with sustained AV exposure remains unresolved due to the absence of longitudinal field studies. Resolving these tensions will require longitudinal field studies in diverse cultural contexts. Priority should be given to addressing Gap [0] concerning cross-cultural threshold transfer and Gap [1] concerning longitudinal behavioral adaptation in AV environments — gaps that carry direct implications for both road safety policy and the design of pedestrian-accepting automated vehicle systems.

\end{document}
% BIB START
@article{02_2003_brosseau_trf,
  author={Brosseau, [surname only]},
  title={Unsignalized midblock crosswalks are the most preferred pedestrian crossing faci},
  year={2003},
  journal={preprint}
}

@article{03_2006_yang_tra,
  author={Yang, [surname only]},
  title={The developed model accurately represents the high rate of pedestrian red light },
  year={2006},
  journal={preprint}
}

@article{06_2013_brosseau_trf,
  author={Brosseau, [surname only]},
  title={The model successfully represents the high rate of pedestrian red light running },
  year={2013},
  journal={preprint}
}

@article{09_2018_jiang_aap,
  author={Jiang, [surname only]},
  title={Text distraction causes the greatest impairment to pedestrian crossing performan},
  year={2018},
  journal={preprint}
}

@article{15_2014_study_trf,
  author={Study, [surname only]},
  title={The average pedestrian crossing speed at signalized intersections in Mumbai is 1},
  year={2014},
  journal={preprint}
}